\scanpage{38}
D'autre part, on aura ici
\[
\beta\in Q(\mathbf R)\subset P'(\mathbf R)
\]
donc
\[
\mathbf R(\unsure{\gamma|\sigma})=\mathbf R(\beta,\gamma)=\mathbf R(\beta,\alpha)=\mathbf R(\alpha)
\]
et par suite
\[
\mathbf R(\unsure{\gamma|\sigma})=\mathbf R(\alpha)
\]
qui est donc $=\mathbf C$ sauf dans le cas exceptionnel précisé plus haut.

\uline{Remarque} Soit $M$ un motif sur $\mathbf R$, correspondant à un \deleted{$G$-module} \deleted{\unsure{associé}} $T_B(M)$. Alors on ne peut avoir $f_{\infty M}=\operatorname{id}_M$ que si $M$ est une somme de puissances tensorielles de $\mathbf Q(2)$.

\figure{\begin{tikzpicture}\draw (0,0)--(.8,.03);\end{tikzpicture}}

\deleted{La relation
\[
f_\infty j(\lambda)f_\infty^{-1}=j(\lambda')
\]
implique, faisant $\lambda=C\in S(\mathbf R)$
\[
f_\infty C f_\infty^{-1}=C^{-1}=\eta C\qquad(\eta=i(-1))
\]
donc
\[
f_\infty\in\underline{\operatorname{Norm}}_{G_{\mathbf R}}(C,C^{-1})
\]
\deleted{Or, \unsure{comme} $K'=\underline{\operatorname{Norm}}_{G_{\mathbf R}}(C,C^{-1})\supset\underline{\operatorname{Norm}}_{G_{\mathbf R}}(C)=K$ qui est un sous-groupe d'indice 1 ou 2. On sait que $K(\mathbf R)$ est un compact maximal dans $G(\mathbf R)$} Considérons $G/i(G_m)=\dot G$ $[=\operatorname{Ker}\varepsilon/\underline i(\mu_2)]$, et les images $\dot f_\infty$, $\dot C$ de $f_\infty$, $C$, donc on a
\[
\dot f_\infty\in\underline{\operatorname{Norm}}_{\dot G_{\mathbf R}}(\dot C,\dot C^{-1})\overset{\mathrm{dfn}}=K'
\]
Or \deleted{on sait que} \unsure{soit}
\[
K=\underline{\operatorname{Norm}}_{\dot G_{\mathbf R}}(\dot C)\subset K'
\]
}
\mspage{37}
